% ============================================================
% 五、实现与实验 (Implementation & Experiments)
% ============================================================
\section{实现与实验}

\subsection{实现环境}

\begin{itemize}
    \item \textbf{编程语言}：Python 3.x（推荐）/ C++ / Julia / MATLAB
    \item \textbf{关键依赖}：NumPy, SciPy, Matplotlib, PyTorch（如适用）
    \item \textbf{运行环境}：macOS / Linux，CPU / GPU（说明型号）
    \item \textbf{代码仓库}：【GitHub 链接，建议公开仓库并附 README】
\end{itemize}

\subsection{实验设置}

\subsubsection{数据集}

描述所使用的数据集（真实数据或合成数据），包括规模、维度、生成方式。

\subsubsection{评估指标}

列出并解释用于衡量方法性能的定量指标（例如：目标函数值 vs 迭代次数、重构误差、分类准确率、运行时间等）。

\subsubsection{对比方法}

如果进行了对比实验，列出所有对比方法及其简要说明。

\subsection{实验结果}

% === 至少一张图 ===

\begin{figure}[htbp]
    \centering
    % \includegraphics[width=0.75\textwidth]{figures/convergence.pdf}
    \caption{【图 1 标题：例如，目标函数值随迭代次数的收敛曲线】}
    \label{fig:convergence}
\end{figure}

说明：图~\ref{fig:convergence} 展示了……

% === 至少一张表 ===

\begin{table}[htbp]
    \centering
    \caption{【表 1 标题：例如，不同方法的性能对比】}
    \label{tab:results}
    \begin{tabular}{lccc}
        \toprule
        方法 & 指标 1 & 指标 2 & 运行时间 (s) \\
        \midrule
        本文方法 & --- & --- & --- \\
        对比方法 A & --- & --- & --- \\
        对比方法 B & --- & --- & --- \\
        \bottomrule
    \end{tabular}
\end{table}

说明：表~\ref{tab:results} 显示了……

\subsection{结果分析与讨论}

\begin{itemize}
    \item 方法是否按理论预期工作？如不符合，可能的原因是什么？
    \item 参数变化（学习率、正则化系数等）对结果有什么影响？
    \item 方法的成功案例与失败案例分别是什么样的？（建议展示具体例子）
    \item 是否存在理论未覆盖的边界情况？
\end{itemize}
